%!TEX root = ../template.tex %
% !TeX spellcheck = e n_US
%%%%%%%%%%%%%%%%%%%%%%%%%%%%%%%%%%%%%%%%%%%%%%%%%%%%%%%%%%%%%%%%%%%%
%% appendix4.tex
%% NOVA thesis document file
%%
%% Chapter with example of appendix with a short dummy text
%%%%%%%%%%%%%%%%%%%%%%%%%%%%%%%%%%%%%%%%%%%%%%%%%%%%%%%%%%%%%%%%%%%%

\typeout{NT FILE appendix4.tex}%

\chapter{2D dielectric function and optical conductivity equivalence}
\label{app:2D_optical_conductivity}

In this appendix, we relate the dielectric function with the optical conductivity in \gls{2d} systems.
The derivation presented here is adapted from unpublished notes by N.~M.~R.~Peres.

Let us consider the Poisson equation

\begin{equation}
    \label{eq:Poisson_equation}
    \nabla  \vdot  \bE(\br,t)=\frac{e}{\epsz}\delta(\br)+\frac{e}{\epsz}n_{\BD}(\br,t)\delta(z),
\end{equation}

\noindent and the continuity equation

\begin{equation}
    -e\frac{\partial n_{\BD}}{\partial t}+\nabla  \vdot  \bJ_{\BD}=0\,,
\end{equation}

\noindent that in the frequency--momentum domain is equivalent to

\begin{equation}
    \omega e n_{\BD}(\bk,\omega)=-\bk  \vdot  \bJ_{\BD}(\bk,\omega)\,.
\end{equation}

\noindent If we invoke Ohm's law $\bJ_{\BD}(\bk,\omega)=\rttensor{\sigma}(\bk,\omega)\bE(\bk,\omega)$ taking into account the tensor nature of the conductivity, and write the electric field in terms of a potential $\phi$ as $\bE(\bk,\omega)=-\ii\bk\,\phi(\bk,z=0,\omega)$, such that

\begin{equation}
    e n_{\BD}(\bk,\omega)=\frac{\ii}{\omega}\bk \vdot \rttensor{\sigma}(\bk,\omega) \vdot \bk\,\phi(\bk,z=0,\omega)\,,
\end{equation}

\noindent Then, we perform a Fourier transform on Eq.~\eqref{eq:Poisson_equation}, while separating the in-plane wavevector from the off-plane one, that we shall denote by $\bk$ and $\kappa_{z}$, respectively.
Doing so, we arrive at

\begin{equation}
    \begin{gathered}
        -k^{2}\phi(\bk,\kappa_{z},\omega)-\kappa_{z}^{2}\,\phi(\bk,\kappa_{z},\omega)=\frac{e}{\epsz}+\frac{\ii}{\epsz\omega} \bk \vdot \rttensor{\sigma}(\bk,\omega) \vdot \bk\, \phi(\bk,z=0,\omega)\Leftrightarrow\\
        \Leftrightarrow\phi(\bk,\kappa_{z},\omega)=-\frac{1}{k^{2}+\kappa_{z}^{2}}\left[\frac{e}{\epsz}+\frac{\ii}{\epsz\omega} \bk \vdot \rttensor{\sigma}(\bk,\omega) \vdot \bk\, \phi(\bk,z=0,\omega)\right].
    \end{gathered}
\end{equation}

\noindent In the previous expression, if we do now the inverse Fourier transform only in the off-plane coordinate, we get

\begin{equation}
    \phi(\bk,z,\omega)=\frac{e}{2\epsz k}\mathrm{e}^{-kz}-\ii\frac{\bk \vdot \rttensor{\sigma}(\bk,\omega) \vdot \bk}{\epsz\omega}\phi(\bk,z=0,\omega)\frac{\mathrm{e}^{-kz}}{2k}\,,
\end{equation}

\noindent which at $z=0$ implies

\begin{equation}
    \begin{gathered}
        \phi(\bk,z=0,\omega)\left(1+\ii\frac{\bk \vdot \rttensor{\sigma}(\bk,\omega) \vdot \bk}{2 k \epsz\omega}\right)=\frac{e}{2 \epsz k}\Leftrightarrow\\
        \Leftrightarrow \phi(\bk,z=0,\omega)=\frac{v_{c}(k)}{\varepsilon(\bk,\omega)}\,,
    \end{gathered}
\end{equation}

\noindent so that we can identify the dielectric function as

\begin{equation}
    \varepsilon(\bk,\omega)=1+\ii\frac{\bk \vdot \rttensor{\sigma}(\bk,\omega) \vdot \bk}{2k\epsz\omega}=1+\ii\frac{k \sigma(\bk,\omega)}{2\epsz\omega}\,,
\end{equation}

\noindent where we considered $\bk \vdot \rttensor{\sigma}(\bk,\omega) \vdot \bk = k^2 \sigma(\bk,\omega)$ to write the second equality.
This expression contrasts with the expression $\varepsilon(\bk,\omega)=1+\ii \frac{\sigma(\bk,\omega)}{(\epsz\omega)}$ that we encountered in the first chapter.
Therein, we have implicitly assumed that the material was \gls{3d}, explaining the discrepancy between the two expressions.
